Deep learning  models are becoming widely used in high-risk settings such as healthcare, transportation, and finance. In these settings, it is important that a model produces calibrated uncertainty to reflect its own confidence and to assist decision making. Confidence regions are a common approach to quantify forecast uncertainty \cite{khosravi2011comprehensive}. A $(1-\alpha)$-confidence region $\Gamma^{1-\alpha}$ for a random variable $y$ is \textit{valid} if it contains $y$'s true value with high probability:  $\Prob(y \in \Gamma^{1-\alpha}) \geq 1-\alpha$. Note one can make  $\Gamma^{1-\alpha}$ infinitely wide to satisfy validity. For the confidence region to be useful, we want to minimize its area while remaining valid; this is known as the \textit{efficiency} of the region. 

Conformal prediction is a powerful method that produces confidence regions with finite-sample guarantees of validity  \cite{vovk2005algorithmic, lei2018distribution}. Furthermore, it makes no assumptions about the forecast model or the underlying data distribution. Its generality, simplicity, and statistical guarantees have made conformal prediction popular for many real world applications including drug discovery \cite{eklund2015application} and robotics \cite{luo2021sample}.

In this paper we present a more calibrated and efficient conformal prediction algorithm for multi-step times series forecasting. This type of problem is ubiquitous - examples include predicting hurricane paths, vehicle trajectories, and financial or epidemic forecasts. To quantify uncertainty, we want a "cone of uncertainty" that covers the entirety of the forecast horizon. \cite{Schaar2021conformaltime} presented an algorithm, CF-RNN, for multi-step time series with the assumption that all time steps are modeled independently. In practice, however, we found that CF-RNN produces confidence regions often too large to be useful, especially when the problem is multivariate or has a long forecast horizon. 

% 1. widely used in economics
% 2. Machine learning for probablisitic time series forecasting, but not related to CP
% 3. There is one work in the CP literature use copula to model dependencies between multiple targets (non temporal) by messoudi et al.

We introduce  \textbf{CopulaCPTS}, a \textbf{Copula}-based \textbf{C}onformal \textbf{P}rediction algorithm for multi-step \textbf{T}ime \textbf{S}eries forecasting. We improve efficiency by utilizing copulas to model the dependency between forecasted time steps. 
A copula is a multivariate cumulative distribution function that models the dependence between multiple random variables. It is widely used in economic forecasting \cite{nelsen2007introduction, patton2012review} and has been introduced to conformal algorithms by \cite{messoudi2021copula} to model correlation between multiple targets in non-temporal settings. Copula processes have been explored in generative models for time series \cite{salinas2019high, drouin2022tactis}. We found that copulas are effective in capturing uncertainty relations between time steps. Our contributions are: 
\begin{itemize}
    \item \textbf{CopulaCPTS} is a general uncertainty quantification algorithm that can be applied to \textit{any} multivariate  multi-step forecaster, with statistical guarantees of validity.
    \item \textbf{CopulaCPTS} produces significantly sharper and more calibrated uncertainty estimates than state-of-the-art baselines on 4 benchmark datasets, both synthetic and real.
    \item We present a re-calibration technique such that \textbf{CopulaCPTS} can produce valid confidence intervals for time series forecasts of varying lengths. 
\end{itemize}

% and is explored in the ML multi-variate time series literature \cite{salinas2019high, drouin2022tactis}. 

% Conformal prediction (CP) is a method for evaluating uncertainty for machine learning predictions and producing prediction sets (for classification) or confidence intervals (for regression) in a distribution-free manner. Conformal prediction, originally introduced by \cite{vovk2005algorithmic}, provides a framework to quantify prediction uncertainty and generate confidence regions without making any assumptions on the data distribution. It has gained popularity because of its finite-sample coverage guarantees and implementation simplicity. Many has worked on extending CP to different tasks and prediction algorithms, such as for the inductive setting \cite{vovk2005algorithmic}, multi-label classification \cite{wang2015comparison}, and multi-target regression \cite{messoudi2020multitarget}. 